\section{Background}

The development of mRNA-LNP vaccines relies on a complex intracellular delivery cascade that translates a genetic sequence into a functional immunological stimulus. A significant hurdle in the field is the high variability in antigen expression profiles despite uniform dosing. This discrepancy stems from the sensitive dependence of expression kinetics on the physicochemical properties of the lipid carrier and the molecular design of the mRNA transcript. \parencite{pardi2015expressionKinetics}

The problem is abstracted as a series of coupled kinetic hurdles: cellular uptake, endosomal escape---where typically only 1--3\% of the payload successfully reaches the cytosol \parencite{Pei2025EndosomalEscapeLNP}---and the subsequent balance between translation and degradation. Understanding these temporal dynamics is essential because effective immune activation is a threshold-dependent process, requiring specific antigen concentrations and presentation durations to trigger robust B-cell and T-cell responses. This modeling effort provides a mechanistic framework to explore how individual biological parameters shape the sharp versus sustained profiles observed in early experimental results, facilitating the rational design of more potent vaccine platforms.